\chapter{研究基础}
\section{已取得的资料与数据质量分析}
\subsection{收益、标签与共同索引}
日收益9,091,641行，观察范围2015年1月5日至2024年7月31日，全市场日期2329个。有效标签6,116,554行，范围2017年2月6日至2024年7月24日，P10正类669,391、P5正类358,973。标签规模与收益行数不同，既涉及历史回归及完成窗口预热，也涉及未来五日结果完整性，不能将其差值全部归因于单一缺失字段。

\reporttable{已验数据底座及对应统计口径}{tab:foundationdata}{
\begin{tabularx}{\linewidth}{p{3.6cm}p{3.2cm}Y}\toprule
对象&规模或日期&统计口径\\\midrule
日收益&9,091,641行&2015-01-05至2024-07-31\\
全市场日历&2329个日期&不是每只股票的有效收益日\\
有效公司日标签&6,116,554行&2017-02-06至2024-07-24\\
P10正类&669,391&同一有效标签面板，约10.94\%\\
P5正类&358,973&同一面板稳健性标签，不另加公司日\\
业绩预告&162,168条&原事件总量，模型输入另按公开日期查询\\
诉讼仲裁&21,219条&行为描述可用，最新进展不能回填首次日\\\bottomrule
\end{tabularx}}

全量另目录重建标签与旧表解压内容一致。独立NumPy参照抽查20家公司19,476标签，残差、阈值和P10/P5无数值或标签不一致。该参照是抽样，不是第二套算法独立全市场重算。共享四折索引已全量读回，标签、ID、结果结束日及split边界一致，提供了正式模型共享目标的底座。

\subsection{新股东补导的公开时间}
最新原包3,607,689行、76字段，前三工作簿各1,048,575数据行，第四461,964行。外层及工作簿内层CRC通过，输入前后SHA不变；独立计数与原值CSV全量来源读回一致，另用独立XML解析抽查16源行的76字段，原值一致。旧4,036,775行及原包仍保留，新旧源历史逐行对齐尚未完成，不以行数减少宣称去重成功。

报告期范围为2015-01-05至2024-12-31，公开日期却延至2026-07-16。图\ref{fig:yearcounts}按源记录的公开年份统计，2024年包含该年7月24日以后的记录，2025/2026也属于原包；它们不能都进入固定风险样本。研究信息观察窗口内公开记录为3,297,707条，窗口外为309,982条。

\reportfigure{新股东原包记录的公开年份分布}{fig:yearcounts}{\FigYearCounts}{新包全量审计聚合。各年计数合计3,607,689；这是源公开年份而非训练样本年份，2024全年和2025/2026记录需另按2024-07-24截止过滤。}

\begin{longtable}{p{2.4cm}r p{2.4cm}r}
\caption{公开年份图的数值表}\label{tab:yearcounts}\\
\toprule 公开年份&记录数&公开年份&记录数\\\midrule\endfirsthead
\toprule 公开年份&记录数&公开年份&记录数\\\midrule\endhead
2015&200,294&2021&392,168\\
2016&284,746&2022&408,995\\
2017&314,438&2023&430,988\\
2018&333,991&2024&427,881\\
2019&345,614&2025&101,705\\
2020&366,529&2026&340\\\bottomrule
\end{longtable}

2025/2026的晚公开记录并不自动说明原包错误，它们可能描述旧报告期的补充或更正，但也暴露出“按报告期筛选历史表”不足以重建当时信息集的问题。35,192条报告截止日晚于公开日也仅列为诊断，需核业务含义，不能直接认定为错行或提前生效。

\subsection{字段缺失与名单结构}
图\ref{fig:missing}使用同一个全量行数分母，展示A股代码、股东ID、流通股比例和来源缺失。股东ID大量缺失包含自然人，流通股比例缺失也与名单口径有关，因此两类缺失不能直接解释为机构图损坏。主图使用总股本比例，PctTotShr在新源中全量数值有效且位于0—1；本次范围检查不等于独立全量重算股份分母。

\reportfigure{新股东原包的字段缺失比例与数量}{fig:missing}{\FigMissing}{分母均为3,607,689条源记录；计数由全量审计得到。图不表示机构身份错误率或公司日模态覆盖率。}

\reporttable{字段与身份诊断的不同分母}{tab:missingdata}{
\begin{tabularx}{\linewidth}{p{3.6cm}p{3.25cm}Y}\toprule
诊断&数量&分母和解释\\\midrule
A股代码缺失&9,318条&全包记录；不自动回填历史证券\\
GDID缺失&1,732,902条&全包记录，包含自然人等\\
企业属性行缺GDID&1,289条&企业属性1,271,367条，不等于金融机构总数\\
流通股比例缺失&2,021,545条&全包记录，不能据非空认证流通名单\\
信息来源缺失&33条&全包记录，需定位原始来源\\
同GDID多原名称&9,018个ID&观察多对多，不直接合并受益主体\\
同原名称多GDID&69个名称&不能以名字直接创建稳定节点\\\bottomrule
\end{tabularx}}

原始总股东1/1候选全期180,652组，研究信息窗口内165,192组。严格十行及十个唯一排名条件下，全期142,838组、窗口内129,500组满足结构筛选。这个条件没有包括可能合法的并列排名，所以不能当作全部可用完整名单数量。33,597组存在重复排名，后续比较54,743个单位中54,601个持股数原值相同，原报告和序号规则因而比单纯去重更重要。

新包没有UpDt、nobs、ID或DisclosureID，下载的词典HTML仅为官方跳转入口；主要名单枚举已有用户提供的官方页面取值，ShRank并列、ShrSer\_Type分组、GDID稳定范围及历史覆写规则仍待正文说明。结构、身份、公开版本和证券对应必须分开核验，当前\texttt{main\_graph\_ready=false}。

\section{程序和研究积累}
已有标签、四折、来源映射、中文标注、裁决和SFT整批预检工具。来源篡改、成员错误、证据不在正文、未完成裁决等情况会被拒绝。模型原型已经测试缺失回退、修正上界、梯度、去自身和独享机构等性质。77项Python合成功能回归包含数据、模型、标注、映射及学校文稿检查，另有PPT源检查，不能将其总数当作金融效果实验数。

旧11特征先导Logit平均AP为0.221784，LightGBM为0.236579，差约1.48个百分点。这属于既有表格比较，不是本研究M的提升；旧先导的独立验证、时序区间和每日名单口径还需按新协议重做。本报告没有正式Qwen、真实持股图或完整风险模型的实证分数。

\section{实验条件、缺口和取得途径}
当前环境约4CPU、8.8GiB内存，无CUDA，可进行源审计、接口和小样本测试。收益、三因子、预告、诉讼与股东授权资料已有，不把未核价格、市值、换手或财报时点列为必要输入。真实60市场日取窗、事件版本流、日装配器、若干同信息对照和全矩阵调度仍待落实。

人工金标目前为0。中文试标工作簿与私有包已经生成，但还未传到本人本机完成实际标注。下一步由A200/B独立42形成初始反馈，按规范完成其余标注和裁决。原版名单、更正和主体通过新版10组定位，再扩展到至少100组/100主体；受限词典、历史版本和证券有效日期由可授权渠道补证，不要求手工清洗全市场原表。

GPU尚未租用，Qwen与正式风险模型未训练。拟在来源、金标和CPU接口检查满足后，使用实际设备测完整序列显存、吞吐、恢复和I/O。24GiB是起始参考，不是已经保证容纳的条件。没有已经落实的国家重点实验室或企业实验室合作，不虚构实践经历或外部支持。

\section{经费与复现安排}
经费包括资料补证、人工标注、GPU抽取/缓存、风险全矩阵及备份。现有授权数据和CPU优先复用；增量费用归属、GPU报价和总小时待核。预算以实测小时$H$、时费$c$及必要余量计算，人工工时由试标记录估计。学生计划自行租算力，但未提供金额、批复或付款证明，具体经费来源及落实状态仍待确认。

原始资料只读，新加工和预测保留在授权私有目录。公开报告和源码只含方法、图示、聚合事实和自写文稿；原公告、股东原行、样本预测与权重不公开。当前学制为两年专业学位硕士，学院、类别/领域、导师、学号、入学年级和正式开题日期未给出，封面保持待填写。本人实际文献研读、导师审核、实践环节和系统手续仍须按个人培养方案办理，不能由程序或文稿代填完成状态。
